\renewcommand{\bibname}{Quellen und Bandzuordnung}
\begin{thebibliography}{9}
\addcontentsline{toc}{chapter}{Quellen und Bandzuordnung}
\footnotesize
\setlength{\itemsep}{4pt}
\setlength{\parskip}{0pt}
\interlinepenalty=10000

\bibitem{project}
Martin Kunze: \emph{Die Grundlagen der Mathematik}.
Band~11, Kapitel~2--4: Äquivalenzaxiome, Ununterscheidbarkeit,
Klassen und Quotientenmengen; ohne Gleichmächtigkeit und
Quotientenabbildungen. Mengenoperationen aus Band~3.
Format nach den Leitfäden zu Band~3, Band~4, Band~5, Band~28 und Band~37.

\bibitem{oxford}
University of Oxford, Mathematical Institute:
\emph{Introduction to University Mathematics (2025--26)},
Abschnitt~2.2, Definitionen~2.8 und 2.10--2.12 sowie anschließender
Partitionssatz. Grundlage für die Begriffe und Klassensätze in
Kapitel~\ref{b11guide:ch:classes}.
\url{https://courses.maths.ox.ac.uk/mod/page/view.php?id=65662}.

\bibitem{petersenmeyer}
Carsten Lunde Petersen und Daniel Meyer:
\emph{On The Notions of Mating}.
Annales de la Faculté des Sciences de Toulouse 21, Sonderheft~5
(2012), S.~839--876.
Abschnitt~2.1, S.~840: Durchschnitte und erzeugte Äquivalenzen;
Definition~2.1, S.~841: Sättigung und größtes gesättigtes Teilstück.
Vergleichsstellen zu Definition~\ref{b11guide:def:saturation}
und Theoremen~\ref{b11guide:thm:bounds},
\ref{b11guide:thm:intersection}, \ref{b11guide:thm:hull}.
Die topologischen und dynamischen Inhalte werden nicht verwendet.
\url{https://www.numdam.org/item/10.5802/afst.1355.pdf}.

\bibitem{cambridge2009}
University of Cambridge: \emph{Mathematical Tripos Part IA 2009,
List of Questions}, Paper~4, Section~I, Aufgabe~1E, S.~14:
Äquivalenzhülle durch endliche Ketten
(Theorem~\ref{b11guide:thm:chains}).
\url{https://www.maths.cam.ac.uk/undergrad/pastpapers/files/2009/List_IA.pdf}.

\bibitem{cambridge2006}
University of Cambridge: \emph{Mathematical Tripos Part IA 2006,
Paper~4}, Section~II, Aufgabe~5E, S.~3:
Komposition zweier Äquivalenzen und gemeinsam erzeugte Äquivalenz
(Theorem~\ref{b11guide:thm:permuting}).
Die Quelle verwendet die entgegengesetzte Kompositionsreihenfolge.
\url{https://www.maths.cam.ac.uk/undergrad/pastpapers/files/2006/PaperIA_4.pdf}.

\bibitem{mathlib}
The mathlib Community: \emph{Mathlib.Order.Setoid.Basic},
offizielle Dokumentation, abgerufen am 23.~Oktober 2026.
Einträge \path{Setoid.quotientQuotientEquivQuotient}
und \path{Setoid.correspondence}.
Theoreme~\ref{b11guide:thm:iterated}
und~\ref{b11guide:thm:correspondence} verwenden nur die
klassenbezogenen Inhalte, ausgedrückt durch Mengen und Vereinigungen.
Abbildungs- und Ordnungsaussagen werden nicht übernommen.
\url{https://leanprover-community.github.io/mathlib4_docs/Mathlib/Order/Setoid/Basic.html}.

\end{thebibliography}

\begingroup
\footnotesize
\noindent\textbf{Einordnung.}
Die übrigen Kriterien und Rechengesetze sind hier zusammenhängend
formulierte elementare Folgerungen aus den Äquivalenzaxiomen und
den Mengendefinitionen. Die Quellen kontrollieren Begriffe und
zentrale Resultate; nicht jede Formulierung steht dort wörtlich.
Es werden keine neuen Forschungsergebnisse oder historischen
Erstentdeckungen beansprucht.
\endgroup

